\section{Motivation and broader questions}
\label{app:motivation}

\subsection{Observation from within a system}

An embedded observer uses its body, senses, instruments, memory, language, and inference rules to form a description of the world. Can it also determine how those means of representation condition its experience? Its evidence for that question must come through the same means of observation.


Philosophy offers several formulations of this question. Plato distinguishes appearances from what produces them. Aristotle distinguishes kinds of cause \parencite{Plato1992,Aristotle1984}. Husserl examines how the structure of experience enters scientific descriptions \parencite{Husserl1970}. Leibniz relates perspective to sufficient reason and the identity of indiscernibles \parencite{Leibniz1989,LeibnizClarke2000}. Peirce connects the meaning of a concept to its conceivable practical consequences \parencite{Peirce1877,Peirce1878}. Whitehead emphasizes processes and relations \parencite{Whitehead1978}.

The conjecture about self-observation is
\begin{quote}
\emph{No sufficiently expressive representational system can completely justify or perceive the boundary that generates its own representational capacity.}
\end{quote}
The conjecture concerns an observer that must account for its own access. Formal self-reference has related limits in logic and semantics \parencite{Cantor1891,Russell1908,Godel1931,Tarski1936,Lawvere1969,Hofstadter1979,Hofstadter2007}. Kant distinguishes the conditions of experience from objects considered independently of those conditions \parencite{Kant1998}. Second-order cybernetics includes the observer in the system it describes \parencite{Ashby1956,VonFoerster1981,MaturanaVarela1980}. Work on generalized observers and self-referential living systems asks how to formalize that relation \parencite{Elshatlawy2025,Abramsky2025}. A general result for embedded physical observers remains open.

\subsection{Recovery through a specified interface}

A restricted version of the question has an exact answer. Given a source space and a fixed observation map, which differences between source states can an observation distinguish?

Let \(X\) be a source state space, let \(O\) be an accessible observation space, and let \(\pi:X\to O\) be the observation map. When
\[
\pi(x_1)=\pi(x_2)
\qquad\text{for}\qquad
x_1\neq x_2,
\]
the shared reading identifies a fiber containing both states. A function of that reading can return at most one of them as its answer.

The reconstruction theorem assumes a specified interface. A hidden extension shows that the same outputs can arise from different latent models. Whether a particular physical system must have a noninjective self-interface remains a separate question.

\subsection{Competing models and the stopping rule}

Several hypotheses can account for an observer who experiences a limit on its access.
\begin{enumerate}[label=(\alph*)]
\item one foundational observer that generates a contained reality and relinquishes access to its own control;
\item many genuinely independent observers within a shared physical world;
\item a physical universe that produces bounded biological observers;
\item an information-based reality without a singular foundational observer; and
\item a cognitive model through which an embedded mind organizes severe uncertainty and self-reference.
\end{enumerate}
If these models have the same possible reports, those reports alone cannot determine which model produced them.

A comparison requires a difference in predictions.
\begin{quote}
\emph{Treat the single-observer model, intentional forgetting, revelation, and every external-source model as metaphysical hypotheses unless they produce a testable prediction that distinguishes them from competing models, including through different outcome probabilities.}
\end{quote}
Peirce and Popper relate inquiry to practical consequences and the possibility of failure \parencite{Peirce1877,Popper1959}. Preregistered tests of consciousness theories specify outcomes that distinguish the competing predictions \parencite{Cogitate2025,Robinson2025}.

